% =============================================================================
% §1.5 §4 RECAST — "Four Types of Action"   (DRAFT FOR REVIEW — NOT YET APPLIED)
% Generated 2026-05-27. Replaces the old "Three Types of Action" subsubsection
% (lines 127–140 of 1.5 ... -v2.tex: simple appetition / habitual-rational /
% evaluatively complex).
%
% PRINCIPLE: four types carved by the species of desire that SOURCES the action
% (D1; supersedes D6/D7). doxa cross-cuts all four; epithymetic alone may be
% pre-doxastic. poiēsis/praxis terminus axis is orthogonal (→ §6).
%
% INHERITED VERBATIM (verified): DMA 701a29–b1 (T1); DMA 701a7–16 (T4).
% RELOCATED OUT: dark-alley example → §5 (rest-of-PHASE-3); technē-hexis /
%   praxis-hexis dispositional mapping stays in §6 (already there).
% VERIFIED 2026-05-27 against corpus PDFs:
%   BCAP pp.98-99 — all §4 fragments VERBATIM (arousal/lucid p.99; praktón p.98;
%     impossible / under-our-control / éschaton p.99). Register transcription typos
%     (occuppies→occupies, gos→goes) to fix in verbatim_passages.md.
%   Rhetoric I.5 1360b4-6 — "happiness and its constituents" VERBATIM (PDF p.~14).
%   Gross p.26 — print p.26 CONFIRMED, but the quoted line is HEIDEGGER's (block-
%     quoted by Gross, German original supplied beneath), NOT Gross's own prose,
%     and the passage generalizes to "the pathē" with FEAR as its example, NOT
%     anger. T2 below PROVISIONALLY corrected (attribution + generality) — confirm
%     option in chat before apply.
% NE III.2/III.3 — NOT in corpus; quoted fragments + flagged typos (wheter→whether,
%   broughtt→brought) RAISED for your manual verification. Hold apply until cleared.
% Greek macrons in Unicode (ē) and node-notation $A_n$ per target convention;
%   file-wide A1/$A_n$/pathē-pathos normalization is PHASE 5.
% =============================================================================

\subsubsection*{Four Types of Action}

The three-factor schema of \textit{De Anima} III.10 and the \textit{orexis} that moves through it yield, not three, but four structurally distinct types of action, each individuated by the species of desire that sources it. Since $A_4$ is the actualization of desire, the kinds of action follow the kinds of desire: epithymetic from appetite (\textit{epithymia}); thymotic from spiritedness (\textit{thymos}); bouletic from the rational wish for an end (\textit{boulēsis}); prohairetic from the deliberate choice of the means to that end (\textit{prohairesis}). Aristotle's own account already gestures at the division, closing its description of impelled movement by noting that things that desire to act ``make and act sometimes from appetite or impulse and sometimes from wish'' (Aristotle, \textit{De Motu Animalium}, 701a29--b1).

These four are not mutually exclusive compartments but name the desire that \textit{sources} a given action. A single episode may pass through several: feeling pity, I may \textit{wish} to help---a bouletic motion toward the end---and then \textit{reason out} the course that would relieve the suffering---a prohairetic motion toward the means. The agent is never not in a mood, and the general \textit{pathē}---being-affected as such---prepare the organic parts in every action whatever (\textit{De Motu Animalium} 7, 702a17--19). The civic passions are one species of this being-affected, its social kind, each bearing a \textit{toward-whom}; they form no type of their own, for a civic passion motivates a desire whose source then fixes the type---anger sourcing thymotic action, pity sourcing first a bouletic wish and then a prohairetic choice of means. The typology marks the predominant source-desire, not an exclusive one.

A word on \textit{doxa}. \textit{Doxa}---the soul's involuntary taking-as-true---is not a fifth source but cross-cuts all four: any of the four desires may be taken up under a \textit{doxa}. Epithymetic action alone can also fire pre-doxastically---the instinctive flinch, the bare discrimination of the \textit{kritikon}, with no assent given. Yet it can equally run through a \textit{doxa}: judging that this person would be a pleasing companion is an evaluative assent, one that may be true or false. \textit{Thymos}, \textit{boulēsis}, and \textit{prohairesis}, by contrast, are doxastic throughout.

\textbf{Type 1: Epithymetic Action (\textit{epithymia}).} The most compressed configuration of the actualization chain is the one Aristotle describes first. In \textit{De Motu Animalium} 7 he writes:

\begin{adjustwidth}{0.5in}{0in}
And so what we do without reflection, we do quickly. For when a man is actually using perception or imagination or thought in relation to that for the sake of which, what he desires he does at once. For the actualizing of desire is a substitute for inquiry or thinking. I want to drink, says appetite; this is drink, says sense or imagination or thought: straightaway I drink. In this way living creatures are impelled to move and to act, and desire is the last cause of movement, and desire arises through perception or through imagination and thought. And things that desire to act make and act sometimes from appetite or impulse and sometimes from wish (Aristotle, \textit{De Motu Animalium}, 701a29--b1).
\end{adjustwidth}

Three observations follow. First, ``the actualizing of desire is a substitute for inquiry or thinking'': appetitive action does not wait on deliberation, for the \textit{phantasma} that presents the object as drink supplies the functional equivalent of a conclusion, and the desire issues in motion before any evaluatively complex \textit{doxa} can engage. Second, the disjunction ``sense or imagination or thought'' names three channels by which the object may be recognized, but in the epithymetic case any one suffices---perception alone, without deliberation; the recognizing channel varies while the appetite-to-motion structure holds constant. Third, the framing sentence---``what we do without reflection, we do quickly''---presents this unreflective swiftness as the normal case for appetitive action, not as a deviation from a deliberative default.

Epithymetic action has a floor below even the pursuit of the pleasant: the avoidance of the bad in its most immediate, purely sensory form. The animal that startles at a sudden noise, the human being who flinches at an unexpected blow, recoils before any deliberative register is engaged. The genesis of such movement lies in the \textit{kritikon} of perception itself: in the very act of perceiving, in the `now', the discriminative dimension of \textit{aisthēsis} institutes the most basic good and bad, the to-be-pursued and the to-be-avoided (\textit{De Anima} III.7, 431a8--12), and from what it discriminates as bad the \textit{orexis} recoils. \textit{Pathos} here does not generate the movement; it is the very capability of being-moved---the being-affected (\textit{paschein}) that such discrimination and recoil presuppose. This is the one type that can fire pre-doxastically. It need not, however: epithymetic desire equally runs through a \textit{doxa}, as when one judges another to be a pleasing companion and pursues accordingly---an evaluative assent, answerable to how things are, and so capable of being mistaken. What individuates the type is its source in appetite, not the level of doxastic articulation at which it happens to operate.

\textbf{Type 2: Thymotic Action (\textit{thymos}).} Where appetite pursues the pleasant and shuns the painful, \textit{thymos} answers a slight, defends honour, and asserts the self---and it does so always in relation to another. Aristotle's anger arises definitionally from a perceived slight, an appraisal that presupposes the agent's standing, the offender's intention, and the social matrix within which the slight registers as undeserved (Aristotle, \textit{Rhetoric}, 1378a31--b5). \textit{Thymos} is in this sense irreducibly social: its object is disclosed within a \textit{toward-whom}, a Being-with (\textit{Mitsein}) in which others are met as those who can honour or wrong one. % [PROVISIONAL — Gross-rigor correction pending your option choice; see chat]
On Gross's reading the emotions ``tie humans in a unique way to their embodiment\ldots\ by determining the possibilities for moving about a shared world'' (Gross, ``Introduction,'' \textit{Heidegger and Rhetoric}, p.~26); and the Heidegger lecture Gross there translates puts the point in a formula---``the \textit{eidos} of the \textit{pathē} is a disposition toward other humans, a Being-in-the-world'' (Heidegger, quoted in Gross, p.~26). Anger is one such disposition. Because it answers an appraisal, thymotic action is doxastic throughout; and it admits of degree, Heidegger marking the distance between what is grasped ``in a state of arousal, in blind passion'' and what is grasped ``in clear, lucid resolution'' (\textit{BCAP} p.~99).

\textbf{Type 3: Bouletic Action (\textit{boulēsis}).} \textit{Boulēsis} is the rational wish for an end. Its object is the good held as end---``happiness and its constituents,'' which Aristotle names as the end at which every individual man and all men in common aim, and which determines what they choose and what they avoid (Aristotle, \textit{Rhetoric}, I.5, 1360b4--6). What marks \textit{boulēsis} off from \textit{prohairesis} is precisely that it is wish for the \textit{end} and not yet choice of the means: ``we wish to be healthy,'' Aristotle observes, ``but we choose the acts which will make us healthy'' (Aristotle, \textit{Nicomachean Ethics}, III.2, 1111b21--30). Because it reaches toward the end rather than instituting the realizable, \textit{boulēsis} may aim where \textit{prohairesis} cannot: there may be a wish even for impossibles, ``e.g.\ for immortality,'' or for what depends not on us but on others (Aristotle, \textit{Nicomachean Ethics}, III.2, 1111b21--30). Heidegger marks the same boundary---wishing, unlike resolute choice, ``can be directed at something that is impossible'' (\textit{BCAP} p.~99). The bouletic moment is thus the soul's orientation toward an end held as good, whether or not the path to it is yet in view: as when, moved to pity by another's suffering, one wishes the suffering relieved before any means has been weighed.

\textbf{Type 4: Prohairetic Action (\textit{prohairesis}).} \textit{Prohairesis} is the deliberate choice of the means by which a wished end is realized---deliberation completed in a resolution that institutes action. Aristotle's practical syllogism exhibits its form:

\begin{adjustwidth}{0.5in}{0in}
But how is it that thought is sometimes followed by action, sometimes not; sometimes by movement, sometimes not? What happens seems parallel to the case of thinking and inferring about the immovable objects. There the end is the truth seen (for, when one thinks the two propositions, one thinks and puts together the conclusion), but here the two propositions result in a conclusion which is an action---for example, whenever one thinks that every man ought to walk, and that one is a man oneself, straightaway one walks; or that, in this case, no man should walk, one is a man: straightaway one remains at rest. And one so acts in the two cases provided that there is nothing to compel or to prevent (Aristotle, \textit{De Motu Animalium}, 701a7--16).
\end{adjustwidth}

The conclusion here is not a proposition but the action itself---``straightaway one walks.'' What the syllogism stages is the passage from a held end to the realizable act that serves it. Deliberation, Aristotle insists, is not about ends but about what contributes to them: ``having set the end they consider how and by what means it is to be attained'' (Aristotle, \textit{Nicomachean Ethics}, III.3, 1112b12--19), and it ranges only over ``things that are in our power and can be done'' (Aristotle, \textit{Nicomachean Ethics}, III.3, 1112a31). Hence \textit{prohairesis} is bounded in just the way \textit{boulēsis} is not. It goes after the \textit{praktón}, ``that which is decisive for a concern in the moment'' (\textit{BCAP} p.~98); and where wish may reach toward the impossible, choice ``never goes after something that is impossible'' but ``always goes after something that is under our control,'' carrying through to the \textit{éschaton}---the point ``that I grasp, that I genuinely institute through action'' (\textit{BCAP} p.~99). Deliberate choice can also sediment into readiness: the expert pianist no longer weighs each fingering afresh but realizes a passage through means long since chosen and now habitual---prohairetic resolution compressed by practice into fluency. How such sedimentation cultivates a disposition, and how that disposition differs when the end lies in a product rather than in the doing itself, belongs to the bivalence of \textit{hexis} taken up below.

Across all four types the division is by source-desire alone; it is orthogonal to a second distinction---whether the action's end lies beyond it, in a product (\textit{poiēsis}), or is immanent in the doing (\textit{praxis})---which cross-cuts the four and is the matter of the \textit{hexis}-bivalence below. And within each of the doxastic types, what decides whether the desire concretizes into a determinate civic passion---anger, fear, pity, shame---is whether the \textit{doxa} ratifies a content that is evaluatively complex. That condition is what the next subsection makes precise.
